\section{\NEW{Discussion}} \label{sec:discussion}

\begin{anew}
\textbf{New section, as you suggested.} It costs about a page and a half and it
does three things the paper currently cannot: it states the mechanism behind the
headline result, it converts the numbers into something a logistics reader can
act on, and it collects the limitations in one place instead of leaving a referee
to find them. Everything below is drop-in prose; edit freely.
\end{anew}

\subsection{\NEW{Why adaptivity beats information}}

\NEW{The ranking of Section~\ref{sec:basecase} is not a ranking of how much a
method knows. The box robust plan and the budgeted robust plan know only the
support of the multiplier; the two-stage plan and the look-ahead policy both draw
25 scenarios from the same distribution. Yet the two adaptive architectures are
separated from the two committing ones by \prov{ten to twenty} percentage points
of gap, while the two scenario-based methods are separated from each other by
almost nothing on the medium routes. What distinguishes them is \emph{when} the
duration is fixed, not what was known when it was chosen.}

\NEW{The mechanism is visible in the two gap definitions. On the window-free
classes the penalised and duration-only gaps coincide exactly, and the ordering of
the methods compresses; on the tight-window classes the penalty term accounts for
a quarter to two fifths of the gap, and the ordering spreads out. Committing early
is cheap in the rest structure --- the Hours-of-Service pattern of a long route is
largely determined by the corridor, and a plan that fixes it in advance loses
little --- and expensive in the appointments, because a window is a constraint on
an uncertain arrival and the only way to satisfy it robustly is to arrive early
and absorb slack that the realization may not require. Window compliance is
correspondingly the axis on which the methods separate most:
\prov{40.8\%} (RO), \prov{55.5\%} (Greedy), 67.3\% (ROBU), \prov{68.9\%} (LA),
\prov{77.9\%} (SP), against a spread of under three points in Hours-of-Service
violations. A schedule committed before departure survives the travel-time noise
in its rest structure and fails in its appointments.}

\subsection{\NEW{Why the two-stage plan strands}}

\NEW{The two-stage plan's \prov{15.5\%} stranding rate is the most striking number
in Table~\ref{tab:feasibility} and it deserves an explanation, because the obvious
one is wrong. It is not caused by energy noise. Energy enters the model only
through the ECR curve evaluated at the realized leg-average speed, and over the
whole support of the multiplier the per-leg energy deviates by at most $+7\%$
relative to nominal --- the fastest admissible realization, 90\,km/h, consumes
29.4\,kWh on a 25\,km leg against 27.5\,kWh at the nominal 80\,km/h. A bounded
7\% perturbation cannot by itself empty a battery that the plan sized with a
20\% reserve.}

\NEW{What strands the vehicle is the structure of the plan rather than the
realization. The two-stage program commits the charging \emph{pattern} offline and
adapts only durations; when the realization pushes the truck past the point where
the committed pattern is recoverable, lengthening a charge cannot help, because
the missing charge is at a station the plan never selected. The repair MILP exists
precisely for this case, and it fires on \prov{99.2\%} of executions --- which is
itself the diagnosis. A plan that must be repaired on essentially every execution
is operationally a policy, and one whose recourse set is too narrow for the
failure mode it faces.}

\NEW{The same argument, run in the other direction, explains the greedy policy's
mirror-image profile: it never strands (0.0\%) because rule (3) charges whenever
the energy to the next station exceeds the usable state of charge, and it breaches
driving limits on \prov{2.3\%} of runs because it stops only when forced and a
slow leg can exhaust the consecutive-driving budget mid-leg. Greedy is myopically
energy-safe and myopically time-reckless. The two failure modes sit on opposite
tails of the same multiplier, and no policy that ignores one tail is safe against
the other.}

\subsection{\NEW{Operational and policy implications}}

\NEW{Three implications follow for an operator.}

\NEW{\emph{First, adopt a planner before adopting a bigger charger.} Of the
20--23\% electrification penalty an operator would observe when switching a
myopically-scheduled corridor from diesel to battery-electric, roughly half is
the physical cost of charging and roughly half is recoverable by planning
(Section~\ref{sec:diesel}). On a 53\,h medium route the optimal penalty of
\prov{9.4\%} is about \prov{five hours}; the myopic penalty of \prov{22.5\%} is
about twelve. The difference is larger than anything the infrastructure
sensitivity of Section~\ref{sec:sensitivity} can deliver, and it requires no
capital expenditure.}

\NEW{\emph{Second, the right planner is the cheap one.} The look-ahead policy is
the only architecture that is simultaneously near-optimal, reliable and scalable
to the long routes, and Section~\ref{sec:la-config} shows that it can be run at
$|\Xi| = 10$ and $L = 24$\,h for \prov{9\,s} of on-board compute per stop with no
measurable loss. That is a decision every $\approx$40\,min of driving at the
mandated charger spacing, well inside what a cab terminal can absorb while the
truck is already stationary. The offline plans, by contrast, require their full
budget before the vehicle can depart --- up to \prov{four hours} for the budgeted
robust plan on a medium route --- and must be re-incurred whenever the route
changes.}

\NEW{\emph{Third, for the infrastructure planner, power beats density.} At the
spacing already mandated by Regulation (EU) 2023/1804, halving the distance
between charging pools buys \prov{1.2\%} of route duration while tripling the
charger power buys \prov{4.3\%}. The binding constraint on these corridors is how
long a charging stop takes, not how far away the next one is. The return on power
also saturates: 350\,kW captures four fifths of the benefit available at 1\,MW,
and the reason is regulatory as much as physical --- once the charge is shorter
than the mandatory break it runs alongside, further power buys nothing at all,
because the break must be taken regardless. This is a joint HoS/infrastructure
threshold, and it suggests that the marginal euro is better spent bringing
existing sites to 350\,kW than either densifying the network or pushing individual
sites to megawatt class. The caveat is that it displaces the constraint onto grid
connection capacity at each site, which is outside the scope of this model.}

\subsection{\NEW{Limitations}}

\NEW{Six limitations bound the results.}

\NEW{\emph{Daily provisions only.} We model the daily provisions of
Reg.~561/2006 and Dir.~2002/15/EC. The long-route class spans four to five duty
cycles, at which point the weekly rest, the 56\,h weekly driving cap and the 60\,h
working week would begin to bind. The simulator records realized weekly working
time as a diagnostic; \TODO{report the share of long-route schedules that would
breach the 60\,h cap --- it is already logged, and it turns a hole into a measured
limitation.}}

\NEW{\emph{Independent leg multipliers.} Legs are 25\,km and real congestion
events span several consecutive ones. Correlation would hurt precisely the methods
that assume independence --- the budgeted plan's protection level and the
two-stage plan's scenario set --- and would leave the look-ahead policy, which
re-samples at every stop, comparatively unharmed. We therefore expect correlation
to widen rather than narrow the gap we report between adaptive and committing
architectures.}

\NEW{\emph{Leg-average energy model.} Consumption is evaluated at the realized
leg-average speed, so a leg driven at an average of 50\,km/h in stop-and-go
traffic is charged the steady-state consumption at 50\,km/h, which understates it.
Gradient is likewise absent, which matters for a Norwegian corridor in
particular.}

\NEW{\emph{Deterministic charger access.} The access delay $Q_i$ is a known
instance parameter. Charger availability, occupancy and reservation are out of
scope, and \citet{keskin_simulation-based_2021} show that queuing uncertainty
alone can materially affect routing decisions. On a corridor where every truck
seeks the same 45-minute window, this is likely the most consequential omission.}

\NEW{\emph{Makespan objective.} Energy price, driver wage and charging tariffs are
not priced, and time-of-use electricity pricing would in particular give the
committing methods an advantage they do not receive here.}

\NEW{\emph{Fixed sequence.} Routing is given. Extending to joint routing and
scheduling would let the vehicle trade a detour against a better-placed charger,
which is exactly the degree of freedom our fixed-corridor design removes.}

\section{Conclusion} \label{section:conclusion}

\begin{anew}
\textbf{Drop-in conclusion}, to be finalised once the outstanding runs land and
the Norwegian case study is resolved either way.

\medskip
This paper studied the joint scheduling of driver activities and partial
non-linear fast charging for a battery electric truck on a fixed long-haul route
under travel-time uncertainty. We formulated a mixed-integer linear program
capturing the daily provisions of EU Hours-of-Service regulation --- split-break
sequencing, reduced daily rests, extended driving allowances, and the rule that a
sufficiently long charge discharges a mandatory break --- and validated its
regulatory core against an independent exact method, reproducing 39 of 40
published optima and improving on the 40th. Around this model we built five
decision architectures spanning the range from full offline commitment to full
online re-optimisation, evaluated them in a common discrete-event simulator over
\prov{1\,800} instances, and benchmarked them against a perfect-hindsight oracle.

\medskip
Three findings stand out. \emph{Adaptivity dominates information}: the two
policies that re-decide en route stay within \prov{12--15\%} of the hindsight
optimum, against \prov{25\%} for a budgeted robust plan and \prov{35\%} for a box
robust plan, and the cost of committing early is paid almost entirely in customer
time-window compliance rather than in driving time. The rolling-horizon policy is
the only architecture that is simultaneously near-optimal, reliable and scalable
to the longest routes, and it can be deployed at \prov{9\,s} of on-board compute
per stop. \emph{Electrification costs \prov{8--11\%} of makespan at the optimum
but \prov{20--23\%} under myopic scheduling}: mandatory breaks do supply usable
charging windows, and our schedules exploit \prov{68--76\%} of them, but the
access overheads of a fast-charging stop absorb about half of that credit, which
qualifies the more optimistic reading of
\citet{brockmann_break-and-charge_2025}. \emph{For infrastructure, power beats
density}: at the spacing mandated by Regulation (EU) 2023/1804, charger power
outweighs charger density by a factor of roughly four, with returns saturating
once the charge is shorter than the break it runs alongside.

\medskip
Several directions follow. Correlated travel times would test the independence
assumption on which the robust and stochastic architectures rely. Endogenous
charger availability would close the gap between our deterministic access delay
and the queueing behaviour a shared corridor actually produces. A cost objective
incorporating time-of-use electricity pricing would change the trade-off between
charging early and charging cheaply. Finally, relaxing the fixed customer
sequence would let the vehicle trade a routing detour against a better-placed
charging opportunity, which is the degree of freedom this study deliberately
removes.
\end{anew}

\appendix
\section{Constraints linearization} \label{section:appendix1}
The constraints relative to linearization of constraints in the deterministic MILP formulation are summarized here.

\begin{alignat}{2}
    l_i^1 &\leq T^{\mathrm{drv}}_{\mathrm{cons}}\, r_i, &\quad& \, i \in \mathcal{I} \label{eq:l1_ub1}\\
    l_i^1 &\leq c_i^d, &\quad& \, i \in \mathcal{I} \label{eq:l1_ub2}\\
    l_i^1 &\geq c_i^d - T^{\mathrm{drv}}_{\mathrm{cons}}\bigl(1-r_i\bigr), &\quad& \, i \in \mathcal{I} \label{eq:l1_lb}
\end{alignat}
\begin{alignat}{2}
    l_i^2 &\leq T^{\mathrm{drv}}_{\mathrm{sh_2}}\, \rho_i, &\quad& \, i \in \mathcal{I}, \label{eq:l2_ub1}\\
    l_i^2 &\leq s_i^d, &\quad& \, i \in \mathcal{I}, \label{eq:l2_ub2}\\
    l_i^2 &\geq s_i^d - T^{\mathrm{drv}}_{\mathrm{sh_2}}\bigl(1-\rho_i\bigr), &\quad& \, i \in \mathcal{I}. \label{eq:l2_lb}
\end{alignat}
\begin{alignat}{2}
    l_i^4 &\leq T^{spr}_{2}\, \rho_i, &\quad& \, i \in \mathcal{I}, \label{eq:l4_ub1}\\
    l_i^4 &\leq s_i^w, &\quad& \, i \in \mathcal{I}, \label{eq:l4_ub2}\\
    l_i^4 &\geq s_i^w - T^{spr}_{2}\bigl(1-\rho_i\bigr), &\quad& \, i \in \mathcal{I}. \label{eq:l4_lb}
\end{alignat}
\begin{alignat}{2}
    l^5_i &\leq T^{spr}_2\, \rho_i, &\qquad& i \in \mathcal I \label{eq:l5_1}\\
    l^5_i &\leq h_i + o_i, &\qquad& i \in \mathcal I \label{eq:l5_2}\\
    l^5_i &\geq h_i + o_i - T^{spr}_2\,(1 - \rho_i), &\qquad& i \in \mathcal I \label{eq:l5_3}
\end{alignat}

\begin{acmt}
There is no $l^3$: the numbering jumps from $l^2$ to $l^4$ because a third
auxiliary was removed at some point. Harmless, but a referee will notice.
Renumber to $l^1 \ldots l^4$.
\end{acmt}

\section{Valid inequalities for the Rolling-horizon subproblems} \label{sec:valid_inequ}

Constraints \eqref{eq:vi-2} fix the \DEL{upper}\FIX{lower} bound on the number of charging stop on each subsection of the path to what is required due to battery constraints. Constraints \eqref{eq:vi-3}--\eqref{eq:vi-4} ensure that each subsection of the route has the minimum required number of break or rest based on driving time. \FIX{Constraint~\eqref{eq:vi-1} links the charge increment to the charging binary, forbidding any charge at an inactive station.}

\begin{equation}
    e^d_i-e^a_i \leq (E^{cap}-E^{min})y_i,
    \qquad \, i \in \mathcal{K}.
    \label{eq:vi-1}
\end{equation}
\begin{equation}
    \sum_{l \in \mathcal K, l<i} y_l \geq \lceil \frac{1}{E^{cap}-E^{min}}(\sum_{l < i} E_l - (E^{0}-E^{min})) \rceil,
    \qquad \, i \in \mathcal{K}.
    \label{eq:vi-2}
\end{equation}
\begin{equation}
    \sum_{l<i} \rho_l \geq \lceil \frac{\sum_{l<i} D_l}{T^{drv}_{sh2}} \rceil-1,
    \qquad \, i \in \mathcal{I}.
    \label{eq:vi-3}
\end{equation}
\begin{equation}
    \sum_{l<i} r_l \geq \lceil \frac{\sum_{l<i} D_l}{T^{drv}_{cons}} \rceil -1,
    \qquad \, i \in \mathcal{I}.
    \label{eq:vi-4}
\end{equation}

\begin{afix}
\eqref{eq:vi-2} is a \emph{lower} bound --- it says at least so many charging
stops are needed --- and \eqref{eq:vi-1} was stated but never described.
\end{afix}

\begin{atodo}
Report what these are worth. They exist to tighten the relaxation that scores
actions in Section~\ref{sec:RHLA}, and the code comments record that
\eqref{eq:vi-3} is not LP-implied: without it ``the relaxation certifies 4 rests
on routes whose optimum needs 5''. One line --- root relaxation bound with and
without, on a representative instance --- turns an undefended appendix into
evidence.
\end{atodo}

\bibliography{paper4}
\end{document}
